\chapter{Өгөгдөл бэлтгэсэн ажил}
\label{ch:methodology}

\section{Хөгжүүлэлтийн орчин (10-р сар, 4-р долоо хоног)}
Python 3.13, pandas, pyarrow, pytest сангуудаар орчин бэлдэж, кодыг өгөгдлийн давхарга (\texttt{fxdata}) ба
скриптүүд гэж хуваасан. Алхам бүрийг нэг командаар дахин ажиллуулах боломжтой, автомат тесттэй байхаар бичсэн.

\section{Ханшийн өгөгдлийн эх сурвалж сонгосон нь}
Хэд хэдэн эх сурвалжийг харьцуулсан (\tref{tab:sources}).

\begin{table}[H]
\centering
\caption{Ханшийн өгөгдлийн эх сурвалжийн харьцуулалт}
\label{tab:sources}
\small
\begin{tabular}{p{2.6cm}p{5.2cm}p{5.4cm}}
\toprule
\textbf{Эх сурвалж} & \textbf{Давуу тал} & \textbf{Сул тал / шийдвэр} \\
\midrule
Dukascopy & 2015-аас хойш бүтэн, үнэгүй, UTC цагтай, тикийн эзлэхүүнтэй & цоорхой, огцом үсрэлттэй; \textbf{үндсэн эх сурвалж} \\
HistData & 1 минутын өгөгдөл, үнэгүй & цагийн бүс нь UTC биш; \textbf{шалгагч эх сурвалж} \\
OANDA & UTC, найдвартай & бүртгэл, токен шаардлагатай; нөөц \\
TrueFX, FXCM, Yahoo & --- & зөвхөн тик, эсвэл хугацаа дутуу; \textbf{тохироогүй} \\
\bottomrule
\end{tabular}
\end{table}

Dukascopy-оос EUR/USD-ийн BID үнийн 1 цагийн лааг \texttt{dukascopy-node} хэрэгслээр татсан. Dukascopy 2026 онд
өгөгдлийн API-гаа өөрчилсөн тул хэрэгслийн хувилбарыг (1.51.1) тогтоож ашигласан. Одоогоор 2015-01-01 -- 2026-11-06
хооронд 73\,878 лаа бий.

\section{Мэдээний календарь ба өгөгдлийн сан (11-р сар, 1-р долоо хоног)}
2007--2025-04 оны мэдээний календарийг ForexFactory-оос цуглуулсан нээлттэй өгөгдлийн сангаас \cite{ffdataset} авсан.
Мөр бүр мэдээний огноо, цаг, валют, нөлөөний түвшин (бага, дунд, өндөр), бодит утга, хүлээлт, өмнөх утгыг агуулна.
Календарь 2025-04-нд дуусдаг тул 2026-11 хүртэлх хэсгийг нээлттэй эхийн скрейпер \cite{ffscraper} ашиглан
ForexFactory-оос татаж нэмсэн (1\,125 USD/EUR мэдээ, үүнээс 538 нь өндөр нөлөөтэй).

Мэдээний \emph{surprise}-ийг бодит утга ба хүлээлтийн зөрүүг тухайн мэдээний түүхэн тархалтад хувааж тооцсон. Мөн
"бодит утга их гарвал валютад сайн/муу" гэсэн тэмдэглэгээг ашиглан surprise-ийг EUR/USD өсөх ($+$) эсвэл буурах
($-$) чиглэлд хөрвүүлсэн.

Өгөгдлийг SQLite өгөгдлийн санд дараах хүснэгтүүдээр хадгална:
\begin{itemize}
    \item \texttt{candles}: лаа (хэрэгсэл, timeframe, эх сурвалж, UTC цаг, OHLC, эзлэхүүн, цэвэр эсэх);
    \item \texttt{data\_quality\_issues}: илэрсэн асуудал бүр, түүнийг хэрхэн шийдсэн;
    \item \texttt{econ\_events}: мэдээ (UTC цаг, валют, нөлөө, утгууд, surprise, цагийн эх сурвалж);
    \item \texttt{candle\_events}: мэдээ ба лааны холбоос.
\end{itemize}
Дараагийн шатанд ашиглах \texttt{model\_runs}, \texttt{predictions} хүснэгтийн бүтцийг мөн тодорхойлсон.

\section{Өгөгдөл цэвэрлэх модуль (11-р сар, 2-р долоо хоног)}
Модуль нь ханшийн өгөгдлөөс алга цаг, огцом үсрэлт (spike), хөдөлгөөнгүй лааг автоматаар илрүүлж, хоёр эх сурвалжийг
цаг бүрээр тулгадаг. Мэдээний календарийн цагийг UTC болгож, мэдээ бүрийг харгалзах лаатай холбоно. Илэрсэн асуудал,
тэдгээрийг хэрхэн шийдсэнийг \tref{tab:quality}-д нэгтгэв.

\begin{table}[H]
\centering
\caption{Илэрсэн асуудал ба шийдэл}
\label{tab:quality}
\small
\begin{tabular}{p{5.2cm}p{7.6cm}}
\toprule
\textbf{Асуудал} & \textbf{Шийдэл} \\
\midrule
Dukascopy ба HistData-ийн зөрүү дунджаар 4.4 пип, АНУ-ын зуны цагийн үед төвлөрсөн &
HistData-ийн цаг Нью-Йоркийн (зуны цагтай) цаг болохыг илрүүлж хөрвүүлсэн. \\
91 огцом үсрэлт & Хоёр эх сурвалжтай тулгах; хоёуланд нь байвал бодит, нэгд нь л байвал алдаа гэж үзэх. \\
Алга цаг, баярын өдрийн цоорхой & Зохиож нөхөхгүй, тэмдэглэж үлдээсэн. \\
Календарийн UTC зөрүү 2023 оноос зуны улиралд буруу &
Файлын зөрүүг үл тоож, хананы цагийг tzdata-аар Asia/Tehran бүсээс дахин хөрвүүлсэн. \\
Өндөр нөлөөтэй мэдээний ~40\% нь цаггүй (00:00) бичигдсэн &
Албан ёсны хуваарь (жишээ нь NFP 08:30 Нью-Йорк), тухайн мэдээний түүхэн цагаас сэргээсэн. \\
Шинэ календарийн цагийн бүс тодорхойгүй &
Тогтмол хуваарьтай мэдээнүүдийн бодит цагтай харьцуулж автоматаар тодорхойлсон. \\
\bottomrule
\end{tabular}
\end{table}

Цаггүй мэдээний асуудал хамгийн аюултай байсан. Ийм мөрийг өмнөх өдрийн орой руу оноовол загвар мэдээний бодит
утгыг 16 цагийн өмнө "мэдэж" байгаа мэт болж, үр дүн хуурамчаар сайн гарна.

\section{Мэдээний нөлөөний шинжилгээний арга (11-р сар, 2-р долоо хоног)}
Мэдээний нөлөө хэдхэн минутад болдог тул \cite{andersen2003micro} 1 цагийн лаа хэт бүдүүн байв. Иймд HistData-ийн
1 минутын өгөгдлийг 5 минутын лаа болгосон. Мэдээ бараг үргэлж 5-д хуваагдах минутад гардаг тул мэдээ лааны яг
эхэнд таарна.

Нэг 5 минутын лаанд гарсан мэдээнүүдийг (жишээ нь NFP, ажилгүйдлийн түвшин, цалин) нэг үйл явдал гэж нэгтгэсэн.
Учир нь ханш тэдгээрт нэг л удаа хариу үзүүлдэг бөгөөд тусад нь тоолбол нэг хөдөлгөөнийг давхар тоолно.

Ханш өдрийн цагаас хамаарч өөрөө өөр хэмжээгээр хөдөлдөг тул мэдээний дараах 60 минутын хүрээг тухайн цагийн ердийн
хүрээтэй харьцуулсан. Ердийн хүрээг сүүлийн 20 өдрийн (ойролцоогоор нэг сар) медианаар тооцсон. Медиан нь өөр том
мэдээтэй өдрүүдэд бага нөлөөлөгддөг бөгөөд 20 өдөр нь тухайн үеийн зах зээлийн нөхцөлийг тусгана. Ингэснээр мэдээ
ханшийг ердийнхөөс хэд дахин их хөдөлгөдгийг өдрийн цагийн нөлөөнөөс салгаж хэмжсэн.